\newtoggle{withbonus} \toggletrue{withbonus} % set true to show bonus points in grade table

% setting underlying course name (Grundlagenfach Mathematik, §-Unterricht, \ldots)
\newcommand{\mycourse}{Grundlagenfach Informatik}
% name of the class
\newcommand{\myclass}{2. Klassen}
% set date
\newcommand{\mydate}{18. Oktober 2024}

\title{\textbf{Bilder bearbeiten in Python}}
\author{Fachschaft Informatik\\
\mycourse, \myclass}
\date{\mydate}
\maketitle
\thispagestyle{headandfoot}		
\clearpage

\section*{Vorbereitung}
\begin{questions}

\question{getting started}

\begin{itemize}
    \item Laden die Dateien
    \begin{itemize}
        \item \pythoninline{aufgaben_bilder.py}
        \item \pythoninline{image_helper.py}
        \item \pythoninline{wohlen_bw.bmp}
        \item \pythoninline{wohlenGray.png}
        \item \pythoninline{wohlen_farbe.jpg}
    \end{itemize}
    von Moodle herunter und speichere sie alle in \textbf{demselben} Ordner auf Deinem Computer ab.
    \item Öffne nun die Datei\pythoninline{aufgaben_bilder.py} und löse die Aufgaben.
\end{itemize}

\section*{Schwarz-Weiss-Bilder}
\question{weisse und schwarze Pixel zählen}

Schreiben Sie ein Python-Programm, welches das Schwarz-Weiss-Bild
\begin{center}
    \pythoninline{wohlen_bw.bmp}
\end{center} als Liste einliest. Das Programm soll die Anzahl der weissen und die Anzahl der schwarzen Pixel des Bildes ausgeben.

Vorgehen:
\begin{itemize}
    \item Lesen Sie das Bild als Pixelliste ein.
    \item Gehen Sie mit einer Schleife durch die Pixelliste und erhöhen Sie einen Zähler, falls eine 1 (beziehungsweise eine 0) gelesen wird.
    \item Geben Sie das Resultat aus.
\end{itemize}

\begin{solution}
\begin{lstlisting}[language=Python,numbers=none]
pixelListe = load_bw_pixels('wohlen_bw.bmp')
anzahlPixelGesamt = len(pixelListe)
anzahlWeissePixel = 0
for pixel in pixelListe:
        anzahlWeissePixel += 1

anzahlSchwarzePixel = anzahlPixelGesamt - anzahlWeissePixel
print('Anzahl weisse Pixel: ' + str(anzahlWeissePixel) + ', Anzahl schwarze Pixel: ' + str(anzahlSchwarzePixel))
\end{lstlisting}
\end{solution}

\question{Pixelwerte invertieren}

Schreiben Sie ein Python-Programm, welches das Schwarz-Weiss-Bild \pythoninline{wohlen_bw.bmp} als Bild und als Liste einliest. Die Aufgabe ist es, das Bild zu invertieren: Alle weissen Pixel sollen schwarz werden und alle schwarzen Pixel sollen weiss werden. Stellen Sie das originale Bild und das invertierte Bild nebeneinander dar.

\begin{solution}
\begin{lstlisting}[language=Python,numbers=none]
import numpy as np
pixelListeBild1 = load_bw_pixels('wohlen_bw.bmp')
Bild1 = load_bw_image('wohlen_bw.bmp')

pixelListeBild2 = []

for pixel in pixelListeBild1:
    if pixel == 0:
        pixelListeBild2.append(1)
    else:
        pixelListeBild2.append(0)

breite, hoehe = Bild1.size

Bild2 = new_bw_image(breite,hoehe,pixelListeBild2)

Image.fromarray(np.hstack((np.array(Bild1),np.array((Bild2))))).show()
\end{lstlisting}
\end{solution}

\question{random image}

Erstelle Sie ein eigenes Bild mit zufälligen Werten (0 oder 1) mit Hilfe des Python-Moduls \pythoninline{random} und lassen Sie dieses Bild anzeigen.

\begin{solution}
\begin{lstlisting}[language=Python,numbers=none]
import random
pixelListeRandom = []

breite = 800
hoehe = 508
anzahlPixel = breite * hoehe

for k in range(anzahlPixel):
    zufallsbit = random.randint(0,1)
    pixelListeRandom.append(zufallsbit)

BildRandom = new_bw_image(breite,hoehe,pixelListeRandom)
BildRandom.show()
\end{lstlisting}
\end{solution}


\section*{Graustufenbilder}

\question{Graustufen invertieren}
Schreiben Sie ein Python-Programm, welches das Graustufenbild \pythoninline{wohlenGray.png} einliest. Das Programm soll die Graustufenwerte invertieren. Beispielsweise soll aus dem hellen (weissen) Pixel mit Wert 255  der dunkle (schwarze) Pixel mit Wert 0 werden. Der dunkle Pixel mit Werte 1 soll zum hellen Pixel mit Wert 254 werden und so weiter. Stellen Sie das originale Bild und das invertierte Bild nebeneinander dar.

\begin{solution}
\begin{lstlisting}[language=Python,numbers=none]
pixelListeBild1 = load_grayscale_pixels('wohlenGray.png')
Bild1 = load_grayscale_image('wohlenGray.png')
pixelListeBild2 = []

for pixel in pixelListeBild1:
    wert = 255 - pixel
    pixelListeBild2.append(wert)

breite, hoehe = Bild1.size

Bild2 = new_grayscale_image(breite,hoehe,pixelListeBild2)
Image.fromarray(np.hstack((np.array(Bild1),np.array((Bild2))))).show()
\end{lstlisting}
\end{solution}

\question{nur 8 verschiedene Graustufen}

Schreiben Sie ein Python-Programm, welches das Graustufenbild \pythoninline{wohlenGray.png} einliest und mit nur 8 verschiedenen (anstelle von 256 verschiednen) Graustufen darstellt. Die Pixel mit Werten in $\lrc{0, 1, 2, \ldots, 31}$ sollen alle durch den Graustufenwert $0$ dargestellt werden, die Pixel mit Werten in $\lrc{32,33,34,\ldots, 63}$ durch den Wert $32$ und so weiter.

\begin{solution}
\begin{lstlisting}[language=Python,numbers=none]
pixelListeBild1 = load_grayscale_pixels('wohlenGray.png')
Bild1 = load_grayscale_image('wohlenGray.png')
pixelListeBild2 = []

for pixel in pixelListeBild1:
    wert = (pixel // 32) * 32
    pixelListeBild2.append(wert)

breite, hoehe = Bild1.size

Bild2 = new_grayscale_image(breite,hoehe,pixelListeBild2)
Image.fromarray(np.hstack((np.array(Bild1),np.array((Bild2))))).show()
\end{lstlisting}
\end{solution}

\section*{anspruchsvolle Aufgaben}

% \question{Gitterstäbe}
%
% Schreiben Sie ein Python-Programm, welches ein Graustufenbild mit \enquote{Gitterstäben} belegt wie recht in \cref{fig:my_label}. Jeder Gitterstab hat eine Breite von zwei Pixeln und zwischen zwei aufeinanderfolgenden Stäben sind in der Breite jeweils 5 Pixel ohne Stäbe.
%
% \begin{figure}[H]
%     \centering
%     \includegraphics[width=0.9\textwidth]{gitter.png}
%     \caption{Gitterstäbe}
%     \label{fig:my_label}
% \end{figure}
%
% \begin{solution}
% \begin{lstlisting}[language=Python,numbers=none]
% pixelListeBild1 = load_grayscale_pixels('wohlenGray.png')
% Bild1 = load_grayscale_image('wohlenGray.png')
% breite, hoehe = Bild1.size
%
% pixelListeBild2 = []
%
% for index in range(len(pixelListeBild1)):
%     wert = (index % breite) % 7
%     if wert == 0 or wert == 1:
%         pixelListeBild2.append(0)
%     else:
%         pixelListeBild2.append(pixelListeBild1[index])
%
% Bild2 = new_grayscale_image(breite,hoehe,pixelListeBild2)
% Image.fromarray(np.hstack((np.array(Bild1),np.array((Bild2))))).show()
% \end{lstlisting}
% \end{solution}


\question{Primzahlen als Pixel}

Schreiben Sie eine Python-Funktion
\begin{lstlisting}[language=Python]
def ist_primzahl(zahl),
\end{lstlisting}
welche für eine gegebene natürliche Zahl \pythoninline{zahl} entscheidet, ob \pythoninline{zahl} eine Primzahl ist (\pythoninline{return True}) oder nicht (\pythoninline{return False}).

\begin{lstlisting}[language=Python,numbers=none]
Testfälle:

ist_primzahl(0) # return False
ist_primzahl(23) # return True
ist_primzahl(97) # return True
ist_primzahl(91) # return False
\end{lstlisting}

Erstellen Sie eine leere Liste. Diese Liste soll am Ende ein Schwarz-Weiss-Bild von $200\times 200$ Pixeln kodieren und somit eine Länge von $40000$ haben. Der $k$-te Eintrag der Liste für
\begin{align*}
    k \in\lrc{0, 1, 2, \ldots, 39999}
\end{align*}
soll $0$ (schwarz) sein, falls $k$ eine Primzahl ist und $1$ (weiss) sonst.

\begin{solution}
\begin{lstlisting}[language=Python,numbers=none]
def ist_primzahl(zahl):
    for k in range(2,zahl):
        if (zahl % k) == 0:
            return False
    return True

breite = 200
hoehe = 200
anzahlPixel = breite * hoehe

pixelListe = []

for index in range(anzahlPixel):
    pixelListe.append(ist_primzahl(index))

Bild = new_bw_image(breite,hoehe,pixelListe)
Bild.show()
\end{lstlisting}
\end{solution}

\section*{Farbige Bilder}
\question{grün verstärken}

Schreiben Sie ein Python-Programm, welches das Farbbild \pythoninline{wohlen_farbe.jpg} und den Grünanteil um 20 Prozent verstärkt.

\begin{solution}
\begin{lstlisting}[language=Python,numbers=none]
pixelListeBild1 = load_rgb_pixels("wohlen_farbe.jpg")
Bild1 = load_rgb_image("wohlen_farbe.jpg")
breite, hoehe = Bild1.size

pixelListeBild2 = []

for (r,g,b) in pixelListeBild1:
    g *= 1.2
    pixel = (int(r), int(g), int(b))
    pixelListeBild2.append(pixel)

Bild2 = new_rgb_image(breite,hoehe,pixelListeBild2)
Image.fromarray(np.hstack((np.array(Bild1),np.array((Bild2))))).show()
\end{lstlisting}    
\end{solution}

\question{Sepia-Filter}

In dieser Aufgabe wollen wir einen Sepia-Filter erstellen und für das farbige Bild \pythoninline{wohlen_farbe.jpg} testen. Wie genau ein Bild durch diesen Filter verändert wird können Sie unter \url{https://de.wikipedia.org/wiki/Sepia_(Fotografie)} nachlesen.

\begin{solution}
\begin{lstlisting}[language=Python,numbers=none]
pixelListeBild1 = load_rgb_pixels('wohlen_farbe.jpg')
Bild1 = load_rgb_image('wohlen_farbe.jpg')
breite, hoehe = Bild1.size

pixelListeBild2 = []

for (r,g,b) in pixelListeBild1:
    r = r * 0.393 + g * 0.769 + b * 0.189
    g = r * 0.349 + g * 0.686 + b * 0.168
    b = r * 0.272 + g * 0.534 + b * 0.131
    pixel = (int(r), int(g), int(b))
    pixelListeBild2.append(pixel)

Bild2 = new_rgb_image(breite,hoehe,pixelListeBild2)
Image.fromarray(np.hstack((np.array(Bild1),np.array((Bild2))))).show()
\end{lstlisting}    
\end{solution}


\end{questions}

